Disclosure appears to resolve this problem without any regulatory intervention. Suppose quality
is verifiable, so that a seller who makes a claim about it can be held to that claim. A seller
of the highest-quality product will disclose, because staying silent means being valued at the
average of the pool rather than at what the product is actually worth. Buyers, anticipating
this, understand that any seller who remains silent is not of the highest quality, and revise
downward what they will pay for silence. That revision gives the second-highest-quality seller
the same incentive, and disclosing becomes worthwhile for them too. The argument repeats. Each
round of disclosure worsens the expected quality of the remaining silent pool, which strengthens
the incentive for the best of those who remain to separate themselves. The process continues
until all sellers but possibly the very worst have disclosed. This is the unraveling result of
\citet{grossman_informational_1981} and \citet{milgrom_good_1981}. Even sellers with bad news
disclose, because they anticipate that silence will be read as worse news still. If it holds,
market forces alone deliver near-complete disclosure and a mandate is redundant.
